	\chapter {Selenium}
	\label{ch:selenium}
      Selenium is a set of different software tools each with a different approach
       to supporting test automation. The entire suite of tools allows many
       options for locating UI elements and comparing expected test results
       against actual application behavior \cite{seleniumSite}.
       Selenium provides implementation both of C\&R \ref{sec:captureReplay}
       \ref{sec:programTests} and programmable tests models.  
       
       \textbf{Selenium IDE} - is a development environment with graphical
       interface for building test scripts.
		Selenium IDE has a recording feature,  which records user actions as they are
		performed and then exports them as a reusable script in one of many programming languages that can be later executed.
		An example of recorded test \ref{lst:capturereplay} opens a URL type two
		numbers to fields with ``number1'' and ``number2'', then -
		clicks and ``add'' button and verify the result see Appendix \ref{appendixA}.
		
       \textbf{Selenium  WebDriver} makes direct calls to the browser using browser's native support for automation. WebDriver represents a web browser, hiding specific browser details, behind the interface see
       figure~\ref{fig:webdriver}. Having a unified interface provides
       multi-browser support and allows same tests to be executed in different
       environments.
       
	  \begin{figure}
	  \centering
  		\includegraphics[width=0.75\textwidth]{webdriver_structure}
  		\caption{Web structure}
  		\label{fig:webdriver}
		\end{figure}
		
    	\textbf{Selenium Remote Control (RC)} is an old version of Selenium
    	WebDriver, currently supported only in maintanance mode.
       
       \textbf{Selenium Grid} allows to execute tests on different machines.
       Selenium Grid is useful for projects with large amount of tests or test
       suites that must be run in multiple environments.
       
		Grid allows to add several physical machines to a test cluster. Grid uses the
		term ``hub'' for a central point where all tests are loaded. Hub is responsible
		for distributing the tests across nodes. ``Node'' is  a remote machine with
		specific configuration which is attached to the hub see
		figure~\ref{fig:selnium_grid}. Nodes are totally separated from each other and
		may have different operating systems and browsers. Hub ``decides''
		for each test suite on which node it should be executed based on test's configuration.
		By default every node starts eleven browsers : five Firefox, five Chrome and
		one Internet Explorer.
		
		\begin{figure}
		\centering
  		\includegraphics[width=0.75\textwidth, center]{selenium_grid}
  		\caption{Selenium Grid structure}
		\label{fig:selnium_grid}
		\end{figure}
		
		Selenium Grid is published as a separate jar file, so to setup a test hub you
		only need to have JRE(Java Runtime Envirenment) installed. For starting hub
		run selenium-server with hub parameter see \ref{lst:starthub}. For starting
		a new node specify a ``webdriver'' parameter and the URL of the hub running
		\ref{lst:startnode}.
		
		\begin{lstlisting} [ label={lst:starthub},language=bash, caption=Start hub ]
		 java -jar	selenium-server-standalone-2.30.0.jar -role hub
		\end{lstlisting}

		\begin{lstlisting} [caption=Start node, label={lst:startnode},language=bash]
		java -jar selenium-server-standalone-2.30.0.jar 
		-role webdriver
		-hub http://http://192.168.1.1:4444/grid/register
		-port 5566
		\end{lstlisting}
		

	To create a test suite for the Grid configuration created above see an
	example \ref{lst:starthub2}. The basic steps are the following:
	\begin{itemize}
	  \item Specify a url of tested web page line 6.
	  \item Specify a url of a hub line 7.
	  \item Specify settings of the test, for example a browser line 8.
	  \item Create an instance of remote webDriver line 9.
	  \item Open a tested web page line 14.
	  \item Compare actual value on a webpage with expected value line 15.
	  \item Close web page line 20.
	\end{itemize}
	
	\lstset{style=a1listing}
	\begin{lstlisting} [caption=Selenium test example,label={lst:starthub2},language=java]
public 	class TestExample throws MalformedURLException {
 	WebDriver driver;
 	 String UIUrl,nodeURL;
 	 @BeforeTest public void setUp() {
 	 	UIUrl="http://app.example.com/hellopage"; 
		hubURL="http://192.168.1.2:5566/wd/hub";
		DesiredCapabilities capability=DesiredCapabilities.firefox();
		driver=new RemoteWebDriver(new URL(hubURL),capability);
	}
		
	@Test
	public void test1() {
		driver.get(UIUrl);
		Assert.assertEquals("Welcome",driver.getTitle());
	}
	
	@AfterTest
	public void afterTest() {
		driver.quit();
	}
}
	\end{lstlisting}

	As it shown above setting up the basic configuration of Selenium Grid is fairly
	easy. After you have opened the tested Web page with ``driver.get(UIUrl)'', you
	can use ``Find Element" or ``Find Elements" methods with ``By" query object for
	locating elements on the page. For example to find an element with a
	class ``profile'' you need to use By.className query object see example
	~\ref{lst:classSearch}. Selenium  supports searching elements by id, tag,
	class, Xpath and all implications from section~\ref{sec:challenges} applies to it. 
	There is no ideal strategy for searching the required element, a developer
	decides which one to choose based on requirements.
	
	\lstset{style=a1listing}
	\begin{lstlisting} [caption=Search element by class,label={lst:classSearch}]
WebElement avatarElement = 
	 driver.findElement(By.className("profile"));
	\end{lstlisting}